% Options for packages loaded elsewhere
\PassOptionsToPackage{unicode}{hyperref}
\PassOptionsToPackage{hyphens}{url}
\PassOptionsToPackage{dvipsnames,svgnames,x11names}{xcolor}
%
\documentclass[
  10pt,
  a4paper,
]{article}
\usepackage{amsmath,amssymb}
\usepackage{iftex}
\ifPDFTeX
  \usepackage[T1]{fontenc}
  \usepackage[utf8]{inputenc}
  \usepackage{textcomp} % provide euro and other symbols
\else % if luatex or xetex
  \usepackage{unicode-math} % this also loads fontspec
  \defaultfontfeatures{Scale=MatchLowercase}
  \defaultfontfeatures[\rmfamily]{Ligatures=TeX,Scale=1}
\fi
\usepackage{lmodern}
\ifPDFTeX\else
  % xetex/luatex font selection
\fi
% Use upquote if available, for straight quotes in verbatim environments
\IfFileExists{upquote.sty}{\usepackage{upquote}}{}
\IfFileExists{microtype.sty}{% use microtype if available
  \usepackage[]{microtype}
  \UseMicrotypeSet[protrusion]{basicmath} % disable protrusion for tt fonts
}{}
\makeatletter
\@ifundefined{KOMAClassName}{% if non-KOMA class
  \IfFileExists{parskip.sty}{%
    \usepackage{parskip}
  }{% else
    \setlength{\parindent}{0pt}
    \setlength{\parskip}{6pt plus 2pt minus 1pt}}
}{% if KOMA class
  \KOMAoptions{parskip=half}}
\makeatother
\usepackage{xcolor}
\usepackage[margin=2.2cm]{geometry}
\usepackage{color}
\usepackage{fancyvrb}
\newcommand{\VerbBar}{|}
\newcommand{\VERB}{\Verb[commandchars=\\\{\}]}
\DefineVerbatimEnvironment{Highlighting}{Verbatim}{commandchars=\\\{\}}
% Add ',fontsize=\small' for more characters per line
\newenvironment{Shaded}{}{}
\newcommand{\AlertTok}[1]{\textcolor[rgb]{1.00,0.00,0.00}{\textbf{#1}}}
\newcommand{\AnnotationTok}[1]{\textcolor[rgb]{0.38,0.63,0.69}{\textbf{\textit{#1}}}}
\newcommand{\AttributeTok}[1]{\textcolor[rgb]{0.49,0.56,0.16}{#1}}
\newcommand{\BaseNTok}[1]{\textcolor[rgb]{0.25,0.63,0.44}{#1}}
\newcommand{\BuiltInTok}[1]{\textcolor[rgb]{0.00,0.50,0.00}{#1}}
\newcommand{\CharTok}[1]{\textcolor[rgb]{0.25,0.44,0.63}{#1}}
\newcommand{\CommentTok}[1]{\textcolor[rgb]{0.38,0.63,0.69}{\textit{#1}}}
\newcommand{\CommentVarTok}[1]{\textcolor[rgb]{0.38,0.63,0.69}{\textbf{\textit{#1}}}}
\newcommand{\ConstantTok}[1]{\textcolor[rgb]{0.53,0.00,0.00}{#1}}
\newcommand{\ControlFlowTok}[1]{\textcolor[rgb]{0.00,0.44,0.13}{\textbf{#1}}}
\newcommand{\DataTypeTok}[1]{\textcolor[rgb]{0.56,0.13,0.00}{#1}}
\newcommand{\DecValTok}[1]{\textcolor[rgb]{0.25,0.63,0.44}{#1}}
\newcommand{\DocumentationTok}[1]{\textcolor[rgb]{0.73,0.13,0.13}{\textit{#1}}}
\newcommand{\ErrorTok}[1]{\textcolor[rgb]{1.00,0.00,0.00}{\textbf{#1}}}
\newcommand{\ExtensionTok}[1]{#1}
\newcommand{\FloatTok}[1]{\textcolor[rgb]{0.25,0.63,0.44}{#1}}
\newcommand{\FunctionTok}[1]{\textcolor[rgb]{0.02,0.16,0.49}{#1}}
\newcommand{\ImportTok}[1]{\textcolor[rgb]{0.00,0.50,0.00}{\textbf{#1}}}
\newcommand{\InformationTok}[1]{\textcolor[rgb]{0.38,0.63,0.69}{\textbf{\textit{#1}}}}
\newcommand{\KeywordTok}[1]{\textcolor[rgb]{0.00,0.44,0.13}{\textbf{#1}}}
\newcommand{\NormalTok}[1]{#1}
\newcommand{\OperatorTok}[1]{\textcolor[rgb]{0.40,0.40,0.40}{#1}}
\newcommand{\OtherTok}[1]{\textcolor[rgb]{0.00,0.44,0.13}{#1}}
\newcommand{\PreprocessorTok}[1]{\textcolor[rgb]{0.74,0.48,0.00}{#1}}
\newcommand{\RegionMarkerTok}[1]{#1}
\newcommand{\SpecialCharTok}[1]{\textcolor[rgb]{0.25,0.44,0.63}{#1}}
\newcommand{\SpecialStringTok}[1]{\textcolor[rgb]{0.73,0.40,0.53}{#1}}
\newcommand{\StringTok}[1]{\textcolor[rgb]{0.25,0.44,0.63}{#1}}
\newcommand{\VariableTok}[1]{\textcolor[rgb]{0.10,0.09,0.49}{#1}}
\newcommand{\VerbatimStringTok}[1]{\textcolor[rgb]{0.25,0.44,0.63}{#1}}
\newcommand{\WarningTok}[1]{\textcolor[rgb]{0.38,0.63,0.69}{\textbf{\textit{#1}}}}
\usepackage{longtable,booktabs,array}
\usepackage{calc} % for calculating minipage widths
% Correct order of tables after \paragraph or \subparagraph
\usepackage{etoolbox}
\makeatletter
\patchcmd\longtable{\par}{\if@noskipsec\mbox{}\fi\par}{}{}
\makeatother
% Allow footnotes in longtable head/foot
\IfFileExists{footnotehyper.sty}{\usepackage{footnotehyper}}{\usepackage{footnote}}
\makesavenoteenv{longtable}
\setlength{\emergencystretch}{3em} % prevent overfull lines
\providecommand{\tightlist}{%
  \setlength{\itemsep}{0pt}\setlength{\parskip}{0pt}}
\setcounter{secnumdepth}{-\maxdimen} % remove section numbering
\ifLuaTeX
  \usepackage{selnolig}  % disable illegal ligatures
\fi
\usepackage{bookmark}
\IfFileExists{xurl.sty}{\usepackage{xurl}}{} % add URL line breaks if available
\urlstyle{same}
\hypersetup{
  colorlinks=true,
  linkcolor={blue},
  filecolor={Maroon},
  citecolor={Blue},
  urlcolor={blue},
  pdfcreator={LaTeX via pandoc}}

\author{}
\date{}

\begin{document}

\section{DDTA BA2 semantic vocabulary candidate -
R2}\label{ddta-ba2-semantic-vocabulary-candidate---r2}

\textbf{Status:} PROVISIONAL CANDIDATE AFTER BA2-T3 / NOT CLOSED / BA2
OPEN

\textbf{Derived by:} BA2-T3 cross-corpus regression

\textbf{Repository baseline reviewed:}
\texttt{4d832bcf90109106d543029cb517be32a6fe7ea7}

\textbf{Structural dependency:}
\texttt{BA2\_PROPOSITION\_STRUCTURE\_CANDIDATE\_R1.md}

\textbf{Closed identity dependency:}
\texttt{BAReferent\ +\ BAProposition}

\subsection{1. Purpose}\label{1-purpose}

R2 is the reduced semantic-vocabulary candidate after replaying the
BA2-T2 seed across the facial-access and order-fulfillment corpora. It
deliberately removes semantics that can be represented more economically
without loss and does not close BA2.

The candidate proposition contract is:

\begin{Shaded}
\begin{Highlighting}[]
\NormalTok{BAProposition}
\NormalTok{|{-} semanticOperatorKey   1}
\NormalTok{|{-} participation         1..*}
\NormalTok{|    |{-} roleKey          1}
\NormalTok{|    \textasciigrave{}{-} term             1}
\NormalTok{|{-} polarity              1       [affirmative by default]}
\NormalTok{\textasciigrave{}{-} scopedModifier        0..*    [only normalized local condition/time scope]}
\end{Highlighting}
\end{Shaded}

\texttt{semanticOperatorKey}, \texttt{roleKey}, \texttt{polarity} and
\texttt{scopedModifier} are vocabulary/structural semantics, not new BAE
identity families.

\subsection{2. Plain-language
interpretation}\label{2-plain-language-interpretation}

The Base Analysis does not need a full systems-modeling language. For
the current evidence it needs only to answer four mechanical questions
about each assertion:

\begin{enumerate}
\def\labelenumi{\arabic{enumi}.}
\tightlist
\item
  \textbf{what kind of claim is this?} (\texttt{semanticOperatorKey});
\item
  \textbf{which project meanings participate, and in what roles?}
  (\texttt{participation});
\item
  \textbf{is the claim affirmed or denied?} (\texttt{polarity});
\item
  \textbf{is there a purely local condition or time scope that changes
  when it applies?} (\texttt{scopedModifier}).
\end{enumerate}

Anything more reusable or independently reviewable is represented as
another \texttt{BAProposition} and/or \texttt{BAReferent}, not hidden in
a modifier mini-language.

\subsection{3. Normative operator-family facet
removed}\label{3-normative-operator-family-facet-removed}

BA2-T2 used \texttt{ACTION}, \texttt{RELATION}, \texttt{CONSTRAINT} and
\texttt{CLASSIFICATION} as a provisional organizing facet.

Cross-corpus regression found no assertion or consumer decision that
required this family value once the exact semantic operator key was
available. The facet therefore does not belong to the normative semantic
lower bound.

A tool may still group registry entries for navigation, but such
grouping is non-authoritative metadata.

\begin{Shaded}
\begin{Highlighting}[]
\NormalTok{operator{-}family facet as normative BA semantics: REJECTED / REMOVED}
\NormalTok{exact semantic operator key: RETAINED}
\end{Highlighting}
\end{Shaded}

\subsection{4. Reduced operator-key
seed}\label{4-reduced-operator-key-seed}

The R2 seed contains thirteen method-neutral operator concepts:

\begin{Shaded}
\begin{Highlighting}[]
\NormalTok{transfer}
\NormalTok{produce}
\NormalTok{create}
\NormalTok{observe}
\NormalTok{transition}
\NormalTok{correlate}
\NormalTok{reference}
\NormalTok{dependOn}
\NormalTok{consumeService}
\NormalTok{realize}
\NormalTok{assignResponsibility}
\NormalTok{constrain}
\NormalTok{classify}
\end{Highlighting}
\end{Shaded}

This remains a falsifiable seed, not an exhaustive universal verb
taxonomy.

\subsubsection{4.1 Merge performed by
BA2-T3}\label{41-merge-performed-by-ba2-t3}

\texttt{ownOrManage} is removed as a separate operator.
Ownership/management is represented as a responsibility/authority kind
under \texttt{assignResponsibility}, with polarity preserving positive
or negative assertions.

Conceptually:

\begin{Shaded}
\begin{Highlighting}[]
\NormalTok{operator: assignResponsibility}
\NormalTok{responsibleParty  {-}\textgreater{} project / party / capability}
\NormalTok{responsibilityScope {-}\textgreater{} transport infrastructure}
\NormalTok{responsibilityKind  {-}\textgreater{} ownership | management | other governed authority kind}
\NormalTok{polarity            {-}\textgreater{} negative}
\end{Highlighting}
\end{Shaded}

This preserves the facial-access fact that the project consumes
transport while not owning/managing it without requiring a separate
operator for every authority mode.

\subsection{5. Why the remaining operator keys
survive}\label{5-why-the-remaining-operator-keys-survive}

\begin{itemize}
\tightlist
\item
  \texttt{transfer} preserves conveyance between source and destination
  and covers source verbs such as deliver/send/submit/emit when they
  carry that meaning.
\item
  \texttt{produce} preserves result/output production and is distinct
  from merely establishing a new identity.
\item
  \texttt{create} preserves establishment of a new project-semantic item
  or event occurrence.
\item
  \texttt{observe} preserves side-effect-free
  reading/querying/observation pressure distinct from state creation.
\item
  \texttt{transition} preserves governed state/lifecycle change.
\item
  \texttt{correlate} preserves same-request/evaluation/operation
  identity matching and remains distinct from generic reference.
\item
  \texttt{reference} preserves explicit pointing without implying
  correlation or dependency.
\item
  \texttt{dependOn} preserves prerequisite/dependency semantics.
\item
  \texttt{consumeService} preserves actual service/capability
  consumption without transferring ownership.
\item
  \texttt{realize} preserves abstract-to-concrete realization such as
  connectivity to current Ethernet/Wi-Fi realization.
\item
  \texttt{assignResponsibility} preserves responsibility/authority
  placement and absorbs ownership/management modes.
\item
  \texttt{constrain} preserves independently queryable restrictions on
  behavior/state/acceptance/applicability.
\item
  \texttt{classify} preserves reusable method-neutral semantic kind
  without creating BA1 subtype families.
\end{itemize}

No reviewed fact forced a fourteenth operator.

\subsection{6. Operator-scoped role
contracts}\label{6-operator-scoped-role-contracts}

BA2-T3 keeps role validity operator-scoped. The table below records the
minimum contract that is sufficient for the reviewed corpora. Maxima are
tightened only where the operator meaning itself requires it; otherwise
future evidence may legitimately widen them.

\begin{longtable}[]{@{}llll@{}}
\toprule\noalign{}
Operator & Required roles & Optional roles & Current cardinality rule \\
\midrule\noalign{}
\endhead
\bottomrule\noalign{}
\endlastfoot
\texttt{transfer} & \texttt{source}, \texttt{destination},
\texttt{content} & none & source \texttt{1}; destination \texttt{1..*};
content \texttt{1..*} \\
\texttt{produce} & \texttt{actor}, \texttt{result} & \texttt{input} &
actor \texttt{1}; result \texttt{1..*}; input \texttt{0..*} \\
\texttt{create} & \texttt{actor}, \texttt{created} & none & actor
\texttt{1}; created \texttt{1..*} \\
\texttt{observe} & \texttt{actor}, \texttt{observed} & \texttt{result} &
actor \texttt{1}; observed \texttt{1..*}; result \texttt{0..*} \\
\texttt{transition} & \texttt{subject}, \texttt{toState} &
\texttt{actor}, \texttt{fromState} & subject \texttt{1}; toState
\texttt{1}; actor \texttt{0..1}; fromState \texttt{0..1} \\
\texttt{correlate} & \texttt{correlatedItem},
\texttt{correlationContext} & none & correlatedItem \texttt{1..*};
correlationContext \texttt{1} \\
\texttt{reference} & \texttt{referencer}, \texttt{referenced} & none &
referencer \texttt{1}; referenced \texttt{1..*} \\
\texttt{dependOn} & \texttt{dependent}, \texttt{prerequisite} & none &
dependent \texttt{1}; prerequisite \texttt{1..*} \\
\texttt{consumeService} & \texttt{consumer}, \texttt{service} &
\texttt{provider} & consumer \texttt{1..*}; service \texttt{1}; provider
\texttt{0..1} \\
\texttt{realize} & \texttt{abstract}, \texttt{realization} & none &
abstract \texttt{1}; realization \texttt{1..*} \\
\texttt{assignResponsibility} & \texttt{responsibleParty},
\texttt{responsibilityScope}, \texttt{responsibilityKind} & none &
responsibleParty \texttt{1..*}; responsibilityScope \texttt{1};
responsibilityKind \texttt{1} \\
\texttt{constrain} & \texttt{constraintTarget}, \texttt{constraintValue}
& none & constraintTarget \texttt{1}; constraintValue \texttt{1..*} \\
\texttt{classify} & \texttt{classifiedReferent}, \texttt{semanticKind} &
none & classifiedReferent \texttt{1}; semanticKind \texttt{1..*} \\
\end{longtable}

Two additions are forced by the regression rather than by lexical
preference:

\begin{itemize}
\tightlist
\item
  \texttt{input} under \texttt{produce}, because
  provider-response/InventoryBalance input must remain visible when a
  governed result is produced;
\item
  \texttt{responsibilityKind} under \texttt{assignResponsibility},
  because responsibility placement, ownership and management must remain
  distinguishable without separate operators.
\end{itemize}

\subsection{7. Semantic key versus lexical
predicate}\label{7-semantic-key-versus-lexical-predicate}

The regression confirms the separation:

\begin{Shaded}
\begin{Highlighting}[]
\NormalTok{source wording}
\NormalTok{    !=}
\NormalTok{BA semantic key}
\NormalTok{    !=}
\NormalTok{display label / synonym}
\end{Highlighting}
\end{Shaded}

Examples:

\begin{itemize}
\tightlist
\item
  \texttt{deliver}, \texttt{send}, \texttt{submit} and \texttt{emit} may
  map to \texttt{transfer} when they express conveyance;
\item
  \texttt{reads}, \texttt{queries} and observation-like derivation
  pressure map to \texttt{observe} when the source state is inspected
  rather than created;
\item
  \texttt{creates} and some \texttt{records} map to \texttt{create} only
  when a new project-semantic item/event is established;
\item
  \texttt{maps} in the payment corpus is not promoted to a \texttt{map}
  operator: it is represented by \texttt{produce} with explicit
  input/result plus a governed constraint on normalization;
\item
  \texttt{release}, \texttt{consume}, \texttt{increase} or
  \texttt{decrease} are not copied mechanically; they map to
  \texttt{transition} and/or \texttt{constrain} according to the
  governed lifecycle/state meaning.
\end{itemize}

Therefore a source author may change a verb without automatically
changing BA operator identity, and the same surface verb may map
differently when its project meaning differs.

\subsection{8. Polarity}\label{8-polarity}

Explicit polarity survives regression.

It avoids proliferating \texttt{notOwn}, \texttt{notAccept},
\texttt{notHandoff}, \texttt{notReserve}, and similar negative operator
keys.

\begin{Shaded}
\begin{Highlighting}[]
\NormalTok{polarity = affirmative | negative}
\end{Highlighting}
\end{Shaded}

Normative prohibitions whose content is itself a reusable rule remain
\texttt{constrain} propositions; polarity does not replace constraint
semantics.

\subsection{9. Modifier reduction}\label{9-modifier-reduction}

BA2-T3 applies the BA2-T2 promotion rule more aggressively.

\subsubsection{9.1 Retained embedded modifier
kinds}\label{91-retained-embedded-modifier-kinds}

Only two normalized local modifier kinds are retained in the R2 lower
bound:

\begin{itemize}
\tightlist
\item
  \texttt{condition} - a proposition-local applicability/guard condition
  that has no independent identity/reuse;
\item
  \texttt{temporalScope} - a proposition-local before/after/during/until
  scope that has no independent identity/reuse.
\end{itemize}

If either becomes reused, separately reviewed or participant-bearing
project meaning, it is promoted to
\texttt{BAReferent}/\texttt{BAProposition} semantics.

\subsubsection{9.2 Removed from the default modifier
vocabulary}\label{92-removed-from-the-default-modifier-vocabulary}

The following BA2-T2 candidate modifier kinds are not retained as
default embedded semantics:

\begin{itemize}
\tightlist
\item
  \texttt{quantification} - participant counts belong in
  role/cardinality contracts; normative occurrence/quantity rules are
  constraints;
\item
  \texttt{completionOutcome} - governed success/failure/completeness
  rules are explicit constraints when analytically relevant;
\item
  \texttt{atomicity} - explicit \texttt{constrain} proposition;
\item
  \texttt{concurrency} - explicit \texttt{constrain} proposition;
\item
  \texttt{idempotency} - explicit \texttt{constrain} proposition.
\end{itemize}

The order corpus repeats atomicity/concurrency/idempotency as
independent obligations, and facial FR-3.4 requires completion/failure
semantics as an independently queryable rule. Hiding them as local
modifiers would weaken analysis selection and provenance.

\subsection{10. Classification
contract}\label{10-classification-contract}

Classification-as-proposition survives regression:

\begin{Shaded}
\begin{Highlighting}[]
\NormalTok{operator: classify}
\NormalTok{classifiedReferent {-}\textgreater{} \textless{}BAReferent\textgreater{}}
\NormalTok{semanticKind       {-}\textgreater{} \textless{}controlled semantic key/value\textgreater{}}
\end{Highlighting}
\end{Shaded}

A fixed exhaustive semantic-kind enumeration is \textbf{not} required by
the current evidence and would risk importing a general systems-modeling
taxonomy.

The contract instead requires:

\begin{enumerate}
\def\labelenumi{\arabic{enumi}.}
\tightlist
\item
  stable controlled semantic-kind keys with normative method-neutral
  definitions;
\item
  provenance/change to attach to the classification proposition in BA3;
\item
  extension only when a consumer would otherwise have to reinterpret raw
  prose;
\item
  no automatic promotion of semantic kinds to BA1 metaclasses.
\end{enumerate}

Current pressure examples include party/person, capability,
information/resource, behavior/event, store and contract distinctions,
but these examples do not constitute a closed universal enumeration.

\subsection{11. Facial-access replay}\label{11-facial-access-replay}

The camera/recognition branch can be represented without raw-prose
reconstruction:

\begin{Shaded}
\begin{Highlighting}[]
\NormalTok{transfer}
\NormalTok{  source      {-}\textgreater{} CameraSubsystem}
\NormalTok{  destination {-}\textgreater{} RecognitionProcessor}
\NormalTok{  content     {-}\textgreater{} RecognitionCapture}

\NormalTok{correlate}
\NormalTok{  correlatedItem     {-}\textgreater{} RecognitionCapture}
\NormalTok{  correlationContext {-}\textgreater{} RecognitionRequest}

\NormalTok{consumeService}
\NormalTok{  consumer {-}\textgreater{} project endpoint context}
\NormalTok{  service  {-}\textgreater{} LocalConnectivity}

\NormalTok{assignResponsibility [negative]}
\NormalTok{  responsibleParty    {-}\textgreater{} project}
\NormalTok{  responsibilityScope {-}\textgreater{} underlying transport}
\NormalTok{  responsibilityKind  {-}\textgreater{} ownership/management}

\NormalTok{realize}
\NormalTok{  abstract    {-}\textgreater{} LocalConnectivity}
\NormalTok{  realization {-}\textgreater{} WiredEthernet}

\NormalTok{constrain}
\NormalTok{  target {-}\textgreater{} delivery behavior}
\NormalTok{  value  {-}\textgreater{} incomplete delivery != successful completion}

\NormalTok{constrain}
\NormalTok{  target {-}\textgreater{} delivery behavior}
\NormalTok{  value  {-}\textgreater{} confidentiality / integrity / authorized provenance obligations}
\end{Highlighting}
\end{Shaded}

No DFD process/data-store/data-flow/trust-boundary type is required in
the shared core.

\subsection{12. Order-fulfillment
replay}\label{12-order-fulfillment-replay}

Representative source verbs normalize without new operator keys:

\begin{itemize}
\tightlist
\item
  checkout creation -\textgreater{} \texttt{create} plus correlation;
\item
  reservation/payment gate dependency -\textgreater{} \texttt{dependOn}
  and \texttt{consumeService};
\item
  availability read -\textgreater{} \texttt{observe}; AvailabilityResult
  -\textgreater{} \texttt{produce};
\item
  provider response normalization -\textgreater{}
  \texttt{produce(input=provider\ outcome,\ result=governed\ result)}
  plus \texttt{constrain};
\item
  reservation/order lifecycle -\textgreater{} \texttt{transition};
\item
  exact-once, atomicity, concurrency and idempotency -\textgreater{}
  \texttt{constrain};
\item
  FulfillmentTask references Order/Reservation -\textgreater{}
  \texttt{reference};
\item
  physical handoff reuse across fulfillment/inventory/payment
  -\textgreater{} reusable behavior/event referent plus
  \texttt{dependOn}/\texttt{correlate}/\texttt{constrain} as applicable.
\end{itemize}

No corpus clause requires copying \texttt{requestsCapture},
\texttt{decreasesReserved}, \texttt{confirmsHandoff}, \texttt{maps} or
another source predicate as a new BA operator.

\subsection{13. Bounded method-consumer
regression}\label{13-bounded-method-consumer-regression}

The earlier STRIDE-oriented camera transfer projection remains
constructible from shared semantics:

\begin{itemize}
\tightlist
\item
  behavior/event classification for the reusable delivery meaning;
\item
  \texttt{transfer} source/destination/content;
\item
  \texttt{correlate} for RecognitionRequest;
\item
  \texttt{consumeService} + responsibility authority semantics for
  external transport responsibility;
\item
  \texttt{realize} for Ethernet;
\item
  \texttt{constrain} for completion and existing security properties.
\end{itemize}

STRIDE categories remain method-owned downstream semantics.

\subsection{14. Remaining BA2 closure
questions}\label{14-remaining-ba2-closure-questions}

The regression leaves a small closure-review set rather than a new
discovery phase:

\begin{enumerate}
\def\labelenumi{\arabic{enumi}.}
\tightlist
\item
  confirm that the thirteen-key seed is the accepted current-scope
  minimum;
\item
  confirm the operator-role matrix and the
  \texttt{ownOrManage\ -\textgreater{}\ assignResponsibility} merge;
\item
  confirm polarity plus only \texttt{condition}/\texttt{temporalScope}
  as embedded modifier lower bound;
\item
  confirm the semantic-kind \textbf{registry contract} without freezing
  a universal kind enumeration;
\item
  confirm that no remaining ambiguity forces BA1 reopen or BA3 leakage.
\end{enumerate}

These are closure questions over the regressed candidate, not
justification for adding vocabulary.

\end{document}
